\section{Gauss 转移定理}\label{sec:galois-1-7}
本节要证明下述定理.
\subsection{Gauss 转移定理}\label{sec:galois-1-7-1}
\begin{theorem*}
若 $A$ 是唯一分解整环，则 $A[X]$ 也是唯一分解整环.
\end{theorem*}
这个定理一个直接而有趣的推论是，多项式代数 $\mathbb C[X_1,\ldots,X_n]$ 都是唯一分解整环.
迈向这个定理的第一步是描述 $A[X]$ 的不可约元.
\subsection{次数为零的不可约元}\label{sec:galois-1-7-2}
\begin{proposition*}
只假设 $A$ 是整环，且 $x\in A$. 则：
\begin{enumerate}[label=\roman*)]
\item $x$ 在 $A$ 中不可约，当且仅当它在 $A[X]$ 中不可约.
\item $Ax$ 是 $A$ 的素理想，当且仅当 $A[X]x$ 是 $A[X]$ 的素理想.
\end{enumerate}
\end{proposition*}
\begin{proof}
i) 只需证明将``不可约''换成``可约''后所得的等价命题. 若 $x$ 在 $A$ 中可约，写成 $x=ab$，其中 $a,b\notin A^\times$，则因 $A$ 是整环，也有 $a,b\notin A[X]^\times$，故 $x$ 在 $A[X]$ 中可约. 反之，若 $x$ 在 $A[X]$ 中可约，写成 $x=fg$，其中 $f,g\notin A[X]^\times$，则由 $A$ 为整环时成立的公式 $\deg(fg)=\deg(f)+\deg(g)$，可知 $f,g\in A$，所以 $x$ 在 $A$ 中可约.

ii) 下面的引理保证 $A[X]/xA[X]\simeq(A/xA)[X]$. 已知 $B$ 为整环蕴含 $B[X]$ 为整环，因此 $A[X]/xA[X]$ 为整环当且仅当 $A/xA$ 为整环，所证即此.
\end{proof}
\begin{lemma*}
设 $A$ 是任意交换环，$I$ 是 $A$ 的理想. 则映射 $f=\sum_n a_nX^n\mapsto\overline f:=\sum_n[a_n\pmod I]X^n$ 诱导一个同构
\begin{align*}
A[X]/IA[X]\xrightarrow{\sim}(A/I)[X].
\end{align*}
\end{lemma*}
\begin{proof}
通过 $A\longrightarrow A/I\longrightarrow(A/I)[X]$ 将 $(A/I)[X]$ 看作 $A$-代数. 于是 $A[X]$ 的泛性质给出将 $X$ 映到 $X$ 的 $A$-代数同态 $\varphi:A[X]\longrightarrow(A/I)[X]$. 它还将 $a$ 映到 $a\pmod I$，因而将 $f$ 映到 $\overline f$，所以这正是命题中的映射. 它的核包含 $I$，故也包含 $I$ 在 $A[X]$ 中生成的理想 $IA[X]$. 商的泛性质于是保证 $\varphi$ 通过同态 $\overline\varphi:A[X]/IA[X]\longrightarrow(A/I)[X]$ 分解.

映射 $f\mapsto\overline f$ 显然满射，故 $\overline\varphi$ 也满射. 为证明 $\overline\varphi$ 单射，须证明 $\varphi$ 的核恰为 $IA[X]$. 若 $f=\sum_n a_nX^n$，则 $\overline f=0$ 当且仅当 $\forall n,\ a_n\in I$；此时 $f\in I+IX+\cdots\subset IA[X]$.
\end{proof}
\begin{corollary*}
设 $A$ 为唯一分解整环. $A[X]$ 中次数为零的不可约元，恰为 $A$ 中的不可约元 $p\in A$. 对这样的 $p$，理想 $pA[X]$ 是素理想，且 $p$-进赋值 $\nu_p:A[X]\longrightarrow\mathbb Z$ 满足：
\begin{align*}
\nu_p(fg)&=\nu_p(f)+\nu_p(g) &&(f,g\in A[X]),\\
\nu_p(f)&=\min\{\nu_p(a_n),\ n\in\mathbb N\} &&\left(f=\sum_n a_nX^n\right).
\end{align*}
\end{corollary*}
\begin{proof}
第一项断言由本分目上面的命题得到. 第一条公式由第\ref{sec:galois-1-5-1}分目第二个引理的 iii) 得到. 对最后的公式，写 $f=p^{\nu_p(f)}f'$，就归结为证明 $\nu_p(f)=0$ 等价于 $\min\{\nu_p(a_n),n\in\mathbb N\}=0$；取逆否命题，这又等价于证明 $p\mid f$ 当且仅当 $p\mid a_i$ 对所有 $i$ 成立. 而 $p\mid f$ 等价于 $\overline f=0$ 在商 $A[X]/pA[X]=(A/p)[X]$ 中成立，上面的引理说明这等价于对所有 $n$ 都有 $p\mid a_n$.
\end{proof}
\begin{remark*}
与第\ref{sec:galois-1-6-3}分目命题的证明相同，赋值 $\nu_p$ 唯一延拓为函数
\begin{align*}
\nu_p:\operatorname{Frac}(A[X])\setminus\{0\}=K(X)^\times\longrightarrow\mathbb Z,
\end{align*}
满足 $\nu_p(fg)=\nu_p(f)+\nu_p(g)$. 此外，对 $f=\sum_n a_nX^n\in K[X]$，仍有公式
\begin{align*}
\nu_p(f)=\min\{\nu_p(a_n),\ n\in\mathbb N\}.
\end{align*}
事实上，选择 $b\in A\setminus\{0\}$ 使 $bf\in A[X]$，便有
\begin{align*}
\nu_p(f)&=\nu_p(bf)-\nu_p(b)\\
&=\min\{\nu_p(ba_n),n\in\mathbb N\}-\nu_p(b)\\
&=\min\{\nu_p(a_n),n\in\mathbb N\}.
\end{align*}
上式与第\ref{sec:galois-1-6-3}分目的命题于是给出
\begin{align*}
\forall f\in K[X]\setminus\{0\},\quad
f\in A[X]\ \Longleftrightarrow\ (\forall p\in A\text{ 不可约},\ \nu_p(f)\geqslant0).
\end{align*}
\end{remark*}
\subsection{次数大于零的不可约元}\label{sec:galois-1-7-3}
设 $A$ 是唯一分解整环.
\begin{definition*}
称 $f\in A[X]$ 是\newterm{galois-primitive-polynomial}{本原多项式}{primitive polynomial}，如果对所有不可约元 $p\in A$，都有 $\nu_p(f)=0$.
\end{definition*}
\begin{proposition*}
设 $f\in A[X]$ 且 $\deg(f)>0$. 则 $f$ 在 $A[X]$ 中不可约，当且仅当 $f$ 本原且在 $K[X]$ 中不可约.
\end{proposition*}
\begin{proof}
设 $f$ 在 $A[X]$ 中不可约. 则它不与任何不可约元 $p\in A$ 等价（因为 $A[X]^\times=A^\times$），所以不被任何这样的 $p$ 整除，因而本原. 设 $f=gh$ 是 $f$ 在 $K[X]$ 中的一个分解. 对每个不可约元 $p\in A$，有 $\nu_p(f)=0=\nu_p(g)+\nu_p(h)$. 选择 $A$ 中不可约元等价类的一组代表 $P$，并令
\footnote{译注：原 PDF 第53页此处第二个乘积的指数印作 $-\nu_p(f)$，按原式保留；结合上文等式，疑应为 $-\nu_p(h)$.}
\begin{align*}
\lambda:=\prod_p p^{\nu_p(g)}=\prod_p p^{-\nu_p(f)}.
\end{align*}
于是得到分解 $f=(\lambda^{-1}g)(\lambda h)$，其中对每个 $p$，都有 $\nu_p(\lambda^{-1}g)=\nu_p(\lambda h)=0$，故由前面的注，$\lambda^{-1}g,\lambda h\in A[X]$. 因为 $f$ 在 $A[X]$ 中不可约，多项式 $\lambda^{-1}g$ 或 $\lambda f$ 中有一个属于 $A[X]^\times=A^\times$.\footnote{译注：原 PDF 此处印作 $\lambda f$，按原文保留；与前面的分解相对照，疑应为 $\lambda h$.} 因此 $g$ 或 $h$ 中有一个属于 $K^\times$，所以 $f$ 确实在 $K[X]$ 中不可约.

反之，设 $f$ 本原且在 $K[X]$ 中不可约，并设 $f=gh$ 是 $A[X]$ 中的一个分解. 这也是 $K[X]$ 中的分解，所以 $g$ 或 $h$ 中有一个属于 $K^\times$. 不妨设为 $g$. 对每个不可约元 $p\in A$，有 $\nu_p(g),\nu_p(h)\geqslant0$，且因 $f$ 本原，有 $\nu_p(g)+\nu_p(h)=\nu_p(f)=0$. 因而对每个 $p$，都有 $\nu_p(g)=0$，于是也有 $\nu_p(g^{-1})=-\nu_p(g)=0$，故 $g^{-1}\in A$. 从而 $g\in A^\times$，所以 $f$ 确实在 $A[X]$ 中不可约.
\end{proof}
\subsection{定理~\ref{sec:galois-1-7-1}的证明}\label{sec:galois-1-7-4}
设 $A$ 为唯一分解整环. 为证明 $A[X]$ 为唯一分解整环，首先要对非零且不可逆的多项式 $f\in A[X]$ 证明不可约元乘积分解的存在性 \hyperlink{galois-factorization-ex}{(Ex)}. 为此，可先选择 $A$ 中不可约元等价类的一组代表 $P$，并作分解 $f=\prod_{p\in P}p^{\nu_p(f)}\cdot f'$，其中 $f'$ 是本原多项式. 这样就归结为在 $f$ 本原的情形证明分解的存在性（此时因 $f$ 不可逆，故其次数大于零）. 对次数归纳即可做到这一点. 若 $f$ 本原且次数为 $1$，它就是不可约的，分解是平凡的. 若 $\deg(f)>0$，则它或者不可约，此时证明结束；或者可写成 $f=gh$，其中 $g,h$ 不可逆. 但因 $f$ 本原，$g,h$ 也都本原，又因它们不可逆，其次数都大于零，从而都小于 $\deg(f)$，故可应用归纳假设.

现在须证明唯一性 \hyperlink{galois-factorization-un}{(Un)}，为此只须证明 $A[X]$ 的每个不可约元都生成素理想. 对来自 $A$ 的零次不可约元，已经验证了这一点. 因此设 $f\in A[X]$ 是次数大于零的不可约元. 已知它在 $K[X]$ 中也不可约. 因为 $K[X]$ 是主理想整环，理想 $fK[X]$ 是素理想（甚至是极大理想），所以商 $K[X]/fK[X]$ 是整环（甚至是域）. 考虑同态 $\varphi:A[X]\longrightarrow K[X]\longrightarrow K[X]/fK[X]$. 因 $\varphi(f)=0$，它通过同态 $\overline\varphi:A[X]/fA[X]\longrightarrow K[X]/fK[X]$ 分解. 来证明 $\overline\varphi$ 单射. 为此须证明 $\operatorname{Ker}(\varphi)=fA[X]$. 设 $g\in\operatorname{Ker}(\varphi)$. 它在 $K[X]/fK[X]$ 中的像为零，故存在 $h\in K[X]$ 使 $g=fh$. 要证明 $h\in A[X]$. 对每个不可约元 $p\in A$，因为 $f$ 本原，有 $\nu_p(f)=0$. 因此 $\nu_p(h)\geqslant\nu_p(g)\geqslant0$，最终得到 $h\in A[X]$. 于是得到单射同态 $A[X]/fA[X]\hookrightarrow K[X]/fK[X]$，它表明 $A[X]/fA[X]$ 为整环，所以 $fA[X]$ 是素理想，所证即此.

